\documentclass[border=10pt]{standalone}
\usepackage{ctex}
\usepackage[european]{circuitikz}
\begin{document}
\begin{circuitikz}
% 1. 逻辑门
\draw (0,6) node[european not port] (not1) {};
\node[above] at (not1.north) {not};

\draw (3,6) node[european and port] (and1) {};
\node[above] at (and1.north) {and};

\draw (6,6) node[european nand port] (nand1) {};
\node[above] at (nand1.north) {nand};

\draw (9,6) node[european or port] (or1) {};
\node[above] at (or1.north) {or};

% 2. 多输入门
\draw (6,4) node[european nand port, number inputs=3] (nand3) {};
\node[above] at (nand3.north) {nand 3-in};

% 3. 触发器 flipflops
\draw (0,2) node[flipflop JK] (jk) {};
\node[above] at (jk.north) {flipflop JK};

\draw (4,2) node[flipflop D] (dff) {};
\node[above] at (dff.north) {flipflop D};

% 4. 开关与灯与电源
\draw (0,-1) to[sinusoidal voltage source, l=$\sim$220V] (0,-3);
\draw (3,-2) node[spdt] (sw) {};
\node[above] at (sw.north) {spdt};

\draw (6,-2) node[lamp] (la) {};
\node[above] at (la.north) {lamp};

\draw (9,-2) to[lamp, l=$Y$] (11,-2);

% dipchip / chip
\draw (0,-7) node[dipchip, num pins=16, hide numbers, no marks] (c16) {};
\node[above] at (c16.north) {dipchip};

\end{circuitikz}
\end{document}
